\documentclass[11pt,reqno]{amsart}
\usepackage{amsmath,amssymb,mathtools}
\usepackage[italian]{babel}
\usepackage{microtype}
\usepackage{tikz}
\usepackage{placeins}
\usetikzlibrary{arrows.meta,positioning,calc,decorations.pathreplacing}
\usepackage[letterpaper,textwidth=6.25in,textheight=8.85in,centering]{geometry}
\usepackage[hidelinks,unicode]{hyperref}
\usepackage{booktabs,array,longtable}
\usepackage{pgfplots}
\pgfplotsset{compat=1.18}
\hypersetup{pdftitle={Pulizia chimica del grafene trasferito con PMMA e PPC},pdfauthor={Codex, su richiesta di Mattia Apicella},pdfsubject={Revisione 2: teoria cinetica, simulazioni condizionali e prove per la sequenza DMF-IPA-THF}}
\numberwithin{equation}{section}
\renewcommand{\abstractname}{Sommario}
\renewcommand{\refname}{Riferimenti}
\renewcommand{\keywordsname}{Parole chiave}
\renewcommand{\datename}{\textit{Data}:}
\newcommand{\SiO}{\ensuremath{\mathrm{SiO}_2}}
\newcommand{\HCl}{\ensuremath{\mathrm{HCl}}}
\newcommand{\degC}{\ensuremath{^{\circ}\mathrm{C}}}
\newcommand{\um}{\ensuremath{\mu\mathrm{m}}}
\newcommand{\nm}{\ensuremath{\mathrm{nm}}}

\title[Grafene: teoria della pulizia a solventi sequenziali]{Pulizia chimica del grafene con residui di PMMA/PPC:\\ teoria della sequenza dei solventi e prove discriminanti}
\author{Codex}
\date{6 ottobre 2026 --- revisione 2}
\thanks{Studio redatto su richiesta di Mattia Apicella per il confronto con Armando Serio. Include calcoli eseguiti e verificati, ma nessun nuovo esperimento sui materiali. Le costanti cinetiche delle simulazioni sono scenari dichiarati, non stime del campione. Questa revisione sostituisce la raccomandazione iniziale di HCl alla luce dei tentativi gi\`a effettuati.}
\keywords{Grafene CVD, PMMA AR-P 672.045, PPC, desorbimento, sequenza dei solventi, bilanci di massa}
\begin{document}
\begin{abstract}
Il problema considerato \`e rimuovere chimicamente residui da grafene
su \SiO/Si, trasferito dal rame con PMMA AR-P 672.045 e PPC.
L'insuccesso riferito di numerosi solventi, acidi e prodotti commerciali
esclude di ripresentarli come nuove soluzioni. L'unico trattamento con
un modesto effetto combina DMF, isopropanolo e THF. Proponiamo di
verificare se il risultato dipenda dalla preparazione di uno stato
polimerico estraibile, dalla sua successiva perdita durante i
risciacqui, oppure da una frazione persistente di diversa natura.
Deriviamo un modello conservativo a tre stati che mostra quando
l'ordine dei liquidi conta e quando l'alternanza aiuta o peggiora.
Un secondo modello separa distacco, trasporto e riadsorbimento;
un bilancio indipendente limita la massa depositabile durante
l'asciugatura. Le simulazioni includono controesempi e controlli
analitici. Ne segue una proposta concreta: confrontare la sequenza
esistente con varianti prive dei passaggi IPA e, se emerge una
mobilizzazione utile, con alternanze a esposizioni complessive uguali.
L'obiettivo \`e amplificare un meccanismo misurabile conservando il
grafene. L'efficacia sul campione reale rimane una domanda sperimentale.
\end{abstract}
\maketitle

\section{La proposta che sceglierei ora}

\textbf{La mia ipotesi di lavoro \`e che la frazione ancora rimovibile
sia limitata dal passaggio tra stati interfaciali del polimero e dalla
successiva estrazione, pi\`u che dalla scelta di un altro solvente
generico.} La prova pi\`u immediata conserva DMF e THF, ma separa
l'effetto dell'IPA intermedio da quello finale. La prova successiva
confronta quattro brevi alternanze con gli stessi liquidi raggruppati
in blocchi, mantenendo uguali esposizioni e ricambi.

Non assumo in partenza che l'IPA sia dannoso. Pu\`o favorire la
precipitazione o il riattacco di materiale gi\`a mobilizzato, oppure
modificare la conformazione di una frazione aderente e renderla
estraibile nel bagno seguente. Queste ipotesi danno previsioni
opposte. Il modello delle Sezioni~\ref{sec:states}--\ref{sec:sim}
le rende quantitative senza attribuire al campione costanti che
nessuno ha misurato.

La modifica da provare per prima, se occorre scegliere un solo
nuovo percorso, \`e \emph{DMF con ricambi freschi, seguito direttamente
da THF, mantenendo inizialmente invariato il risciacquo finale}.
In questo modo si verifica se l'IPA intermedio annulli una
mobilizzazione ottenuta nel DMF. La Tabella~\ref{tab:factorial}
specifica tempi e controlli. Se invece quel passaggio risulta utile,
si passa alla prova di alternanza della Tabella~\ref{tab:cycles}.

La scelta \`e motivata dall'unico percorso che ha gi\`a prodotto un
segnale favorevole. Non \`e deducibile una probabilit\`a di successo
dal suo semplice esistere: il miglioramento riferito non \`e ancora
quantificato e i confronti precedenti non sono descritti abbastanza
da isolare l'effetto della sequenza.

\section{Il caso reale e i risultati negativi da conservare}

Il materiale iniziale indicato \`e PMMA AR-P 672.045, che il produttore
identifica come 950K al 4,5\% in anisolo \cite{Allresist}. Il campione
comprende anche PPC sopra il PMMA, senza contatto intenzionale del
PPC con il grafene. Sono riferiti un passaggio a 90\,\degC\ per
2 minuti, distacco elettrochimico dal Cu e trasferimento su \SiO/Si.
Elettrolita, spessori, eventuali altre cotture e identificazione
chimica delle particelle non sono ancora noti.

Il lavoro di Tyagi, Mi\v{s}eikis, Coletti e collaboratori
\cite{Tyagi2022} corrisponde bene a questo schema PMMA/PPC e ai
passaggi termici indicati. Non attribuiamo automaticamente al caso
attuale tutti i parametri del loro supplemento. Quel lavoro \`e
utile come riferimento di processo; il suo solvente commerciale
non \`e una nuova risposta all'elenco di tentativi di Armando.

Sono gi\`a stati provati, secondo quanto comunicato:
\begin{itemize}
\item HCl per 48 ore e acido acetico glaciale;
\item acetato di etile, cloroformio, DMF, THF, MPK, acetone;
\item IPA/acqua 75/25 e prodotti indicati come AR600.7 e AR300-76,
quest'ultimo sia a temperatura ambiente sia a 80\,\degC.
\end{itemize}
Manteniamo ``MPK'' e ``AR600.7'' come sigle riferite, senza
correggerle arbitrariamente in altri prodotti. Il precedente con
HCl \cite{Xiao2019} resta un risultato pubblicato, ma \textbf{HCl
48 ore viene ritirato come nuova raccomandazione per questo caso}.
La concentrazione mancante limita un confronto quantitativo, senza
cancellare il risultato negativo ricevuto.

L'unica pulizia descritta come debolmente efficace \`e
\begin{equation}\label{eq:real}
\begin{split}
 &\mathrm{DMF}\ (30\,\mathrm{min})
 \longrightarrow \mathrm{IPA}\ (5+3\,\mathrm{min})\\[-2pt]
 &\longrightarrow \mathrm{THF}\ (30\,\mathrm{min})
 \longrightarrow \mathrm{IPA}\ (5+3\,\mathrm{min})
 \longrightarrow \mathrm{N_2}.
\end{split}
\end{equation}
La somma dei bagni \`e 76 minuti. L'agitazione riferita \`e 500 rpm:
non descriviamo quindi il trattamento come un bagno quiescente.
Questa velocit\`a non determina per\`o il moto locale al campione
senza geometria, posizione, volume e tipo di agitatore.

\section{Che cosa aggiunge la letteratura estesa}

\subsection{La sequenza pu\`o cambiare il comportamento del polimero}
Manjkow e collaboratori \cite{Manjkow1987} osservano mediante
ellissometria un passaggio stretto fra dissoluzione e rigonfiamento
del PMMA in miscele solvente/non solvente. Le miscele studiate non
sono quelle di Armando: il lavoro giustifica studiare la traiettoria
di composizione, non applicare qui le sue soglie. Rooks e
collaboratori \cite{Rooks2002} mostrano inoltre perch\'e IPA e
IPA/acqua non si possano classificare con una sola etichetta
indipendente da massa molecolare e stato del resist.

La solubilit\`a del polimero libero non determina il distacco
dall'interfaccia. Simulazioni di desorbimento di catene adsorbite
\cite{Huston2020} trovano regimi cinetici molto diversi al variare
dell'adsorbimento. Si tratta di modelli generici, non di una
parametrizzazione di AR-P 672.045 su grafene. Non \`e corretto
ricavare una barriera moltiplicando l'energia di un monomero per
tutta la catena: distacco progressivo, conformazione e contatti
effettivi contano.

\subsection{Perch\'e il semplice ricambio non \`e gi\`a la soluzione}
Ayodele e collaboratori \cite{Ayodele2022} impiegano PMMA 950K,
distacco elettrochimico dal rame e un estrattore Soxhlet con acetone
distillato. Osservano miglioramenti morfologici, ma il confronto
varia anche temperatura, durata e risciacquo; manca una misura XPS
che certifichi l'assenza del residuo. La storia termica \`e diversa
da quella comunicata per questo caso.

Il risultato del 2025 di Tang e collaboratori \cite{Tang2025}
\`e particolarmente istruttivo: un processo ciclico migliora
l'aspetto macroscopico, ma i rapporti XPS
C=C:C--O:O--C=O sono circa 100:28:8 nel controllo e 100:30:7
nel processo ciclico. I residui rimangono. Questo impedisce di
promettere che solvente rinnovato o assenza di IPA finale bastino
da soli. Il lavoro non testa l'alternanza DMF/IPA/THF qui proposta.

\subsection{Miscela IPA/acqua, etanolo e prestazioni elettriche}
Duan e collaboratori \cite{Duan2022} riportano un miglioramento
con IPA/acqua 75/25, ma dopo acetone e ricottura in vuoto a
200\,\degC. La loro frazione di componenti C--O e C=O resta
misurabile; inoltre il controllo in sola acqua subisce danni.
Il metodo \`e gi\`a nell'elenco di Armando e non viene riproposto.

Merino e collaboratori \cite{Merino2024} studiano residui da
fabbricazione anche con etanolo e THF: documentano miglioramenti
elettrici e chimici parziali, ma anche problemi di distacco con
THF. Non \`e una prova sul nostro PMMA/PPC. Suk e collaboratori
\cite{Suk2013} forniscono un controesempio ancora pi\`u netto:
formammide assorbita nei residui migliora le propriet\`a elettriche
senza corrispondente scomparsa della morfologia residua.
Una migliore curva del transistor, da sola, non decide la pulizia.

\subsection{Acido, identit\`a chimica e diagnosi del residuo}
Nel PMMA l'idrolisi dell'estere laterale non rompe automaticamente
la catena carboniosa. Nei polimetacrilati di Sch\"onemann e
collaboratori \cite{Schoenemann2018}, diversi dal PMMA, la stabilit\`a
all'acido mostra quanto sia rischioso inferire la cinetica dal solo
gruppo funzionale. Anche il PPC ha una risposta all'idrolisi
dipendente dalle condizioni \cite{Jung2006}. Questi studi non
forniscono una nuova miscela di pulizia validata per il campione.

Wang e collaboratori \cite{Wang2017} usano marcatura isotopica del
PMMA per seguirne i residui. Schwartz e collaboratori
\cite{Schwartz2019} identificano localmente polimeri intrappolati
in eterostrutture mediante spettroscopia infrarossa associata
all'AFM. Non studiano il nostro PPC: il loro PC \`e un materiale
diverso. Il punto trasferibile \`e il metodo di identificazione.
Una particella vista in AFM non acquista identit\`a chimica dalla
sola somiglianza con un residuo di PMMA.

\section{Una teoria della sequenza: preparare, estrarre, ribloccare}
\label{sec:states}

\subsection{Una previsione che non richiede conoscere i tassi}
Supponiamo che ogni liquido $j$ rimuova una stessa popolazione di
residuo $m$ con un tasso costante $k_j$, senza memoria:
\begin{equation}\label{eq:scalar}
 \dot m=-k_jm.
\end{equation}
Dopo tempi totali $T_j$ nei vari liquidi,
\begin{equation}\label{eq:orderindependent}
 m_{\mathrm{finale}}=m_0\exp\!\left(-\sum_j k_jT_j\right).
\end{equation}
L'ordine e il numero di alternanze non compaiono. Dunque una
differenza riproducibile fra ordine a blocchi e alternato, a pari
esposizioni, ricambi e condizioni, \emph{falsifica questo modello}.
Non identifica da sola il meccanismo alternativo: concentrazione
del bagno, trasformazioni, stati conformazionali e danno meccanico
possono introdurre memoria. L'osservazione finora riferita non \`e
ancora un confronto controllato di questo tipo.

\subsection{Il modello minimo che permette un effetto dell'ordine}
Dividiamo la massa in tre stati:
$L$, residuo ancora aderente e non prontamente estraibile;
$M$, residuo mobilizzato ma ancora sul campione;
$R$, materiale gi\`a esportato fuori dalla zona di riadesione.
Nel liquido $j$ poniamo
\begin{equation}\label{eq:generator}
 \frac{d}{dt}
 \begin{pmatrix}L\\M\\R\end{pmatrix}
 =K_j\begin{pmatrix}L\\M\\R\end{pmatrix},\qquad
 K_j=\begin{pmatrix}
 -a_j&b_j&0\\a_j&-(b_j+e_j)&0\\0&e_j&0
 \end{pmatrix}.
\end{equation}
Il tasso $a_j$ mobilizza, $b_j$ riblocca, $e_j$ esporta.
Sono non negativi e hanno dimensione tempo$^{-1}$.
La somma $L+M+R$ si conserva e masse inizialmente non negative
rimangono tali. L'assorbimento nello stato $R$ presuppone un
allontanamento efficace: se il polimero resta nel bagno e torna
sulla superficie, occorre aggiungere il comparto della
Sezione~\ref{sec:transport}.

Per due trattamenti D e T, con tempi $t_D,t_T$, l'ordine D poi T
d\`a
\begin{equation}\label{eq:ordered}
 v_{DT}=e^{K_Tt_T}e^{K_Dt_D}v_0,
 \qquad v=(L,M,R)^{\mathsf T}.
\end{equation}
I due esponenziali generalmente non commutano. Per tempi piccoli,
la differenza di ordine parte da
\begin{equation}\label{eq:commutator}
 v_{DT}-v_{TD}=(K_TK_D-K_DK_T)v_0\,t_Dt_T+
 O\bigl((t_D+t_T)^3\bigr).
\end{equation}
Questa \`e una ragione quantitativa per cui due solventi
mediocri separatamente potrebbero essere complementari. Non
assegniamo automaticamente a DMF il ruolo $a$ e a THF il ruolo
$e$: sono proprio i ruoli da mettere alla prova.

\subsection{Lo stesso IPA pu\`o aiutare o ostacolare}
Consideriamo un caso esatto: nel passaggio P non vi \`e esportazione,
ma $L\rightleftarrows M$ con tassi $a,b$; nel successivo G vi sono
$M\to R$ a tasso $k$ e $M\to L$ a tasso $c$. Per tempi $t,u$,
la massa esportata in pi\`u da P poi G rispetto a G poi P \`e
\begin{equation}\label{eq:sign}
 \Delta R=
 \frac{k}{k+c}\bigl(1-e^{-(k+c)u}\bigr)
 \bigl(1-e^{-(a+b)t}\bigr)
 \frac{aL_0-bM_0}{a+b}.
\end{equation}
La formula vale per $a+b>0$ e $k+c>0$, con i casi nulli ottenuti
per continuit\`a. Quando i prefattori sono positivi, il segno
\`e quello di $aL_0-bM_0$.

Se la superficie \`e soprattutto bloccata e P la mobilizza,
il passaggio intermedio aiuta. Se il DMF ha gi\`a prodotto molta
frazione $M$ e l'IPA la riblocca, lo stesso passaggio ostacola
l'estrazione. Questa previsione \`e pi\`u precisa di ``l'IPA fa
precipitare'' o ``l'IPA pulisce''. L'ordine inverso termina in
un liquido diverso: la verifica deve misurare materiale esportato
prima dell'asciugatura o imporre un passaggio finale comune.

\subsection{Perch\'e la contrazione non \`e automaticamente un vantaggio}
Un solvente meno favorevole pu\`o ridurre il volume occupato dalla
catena, ma anche aumentare aggregazione e adesione alla superficie.
Per rendere utile un ciclo occorrono due condizioni: una quota
misurabile passi a uno stato pi\`u estraibile, e il liquido seguente
la allontani prima che torni aderente. Una diminuzione dell'impronta
AFM con aumento dell'altezza e volume conservato suggerisce
riorganizzazione, non ancora rimozione. Il bagno seguente deve
ridurre la massa o il volume con un riscontro chimico.

\section{Calcoli eseguiti e loro conseguenze}
\label{sec:sim}

Abbiamo calcolato gli esponenziali dei generatori con una somma
di matrici stocastiche, detta uniformizzazione, evitando tassi
negativi o perdite artificiali di massa. Il programma allegato
contiene parametri, formule, risultati e verifiche. Nessun
parametro \`e stato adattato a immagini o misure di Armando.

\subsection{Alternanze a tempi totali fissi}
Dividendo in $N$ parti le stesse esposizioni complessive si ottiene
\begin{equation}\label{eq:cycles}
 v_N=\left(e^{K_TT_T/N}e^{K_DT_D/N}\right)^Nv_0.
\end{equation}
Non si confrontano quindi quattro cicli lunghi con un solo ciclo
breve. Per $N\to\infty$,
\begin{equation}\label{eq:trotter}
 v_N\longrightarrow e^{K_TT_T+K_DT_D}v_0.
\end{equation}
Il modello mantiene memoria dello stato mobile fra le fasi.
Un modello che lo azzeri a ogni cambio, anche infinitamente rapido,
avrebbe un limite diverso e non sarebbe equivalente.

La Figura~\ref{fig:cycles} confronta tre scenari adimensionali,
con $T_D=T_T=1$, $v_0=(1,0,0)^{\mathsf T}$ e i tassi della
Tabella~\ref{tab:rates}. Il tempo unitario \`e convenzionale:
non significa un minuto o trenta minuti reali.

\begin{table}[htbp]
\caption{Tassi illustrativi $(a,b,e)$ usati nelle simulazioni.
Non sono propriet\`a misurate di DMF o THF.}
\label{tab:rates}\centering\small
\begin{tabular}{@{}lcc@{}}\toprule
Scenario & Trattamento D & Trattamento T\\\midrule
Complementare &(2; 4; 0{,}02)&(0{,}02; 0{,}1; 4)\\
Preparazione poi estrazione &(2; 0; 0)&(0; 0; 2)\\
Tassi uguali &(0{,}4; 0{,}5; 0{,}3)&(0{,}4; 0{,}5; 0{,}3)\\\bottomrule
\end{tabular}
\end{table}

Nel primo scenario il residuo scende da 0,6680 con una sequenza
a 0,4478 con otto alternanze; a venti risale leggermente a 0,4519.
Nel secondo scenario frammentare la preparazione peggiora la
rimozione. Nel terzo non cambia nulla, come deve accadere con
generatori identici. Questi numeri dimostrano possibilit\`a e limiti
del modello, \emph{non} riduzioni previste per il campione.

Un controesempio \`e anche analitico. Se D mobilizza senza ritorno,
T estrae senza ritorno, e $aT_D=kT_T=w$, la frazione rimossa \`e
\begin{equation}\label{eq:counterexample}
 R_N=1-e^{-w}\left[1+N\bigl(1-e^{-w/N}\bigr)\right].
\end{equation}
Essa diminuisce all'aumentare di $N$: preparare tutta la frazione
prima di estrarla \`e qui pi\`u efficace. Cercare un beneficio dei
cicli senza includere questo caso negativo produrrebbe una teoria
costruita per dare ragione alla proposta.

\begin{figure}[htbp]
\centering
\begin{tikzpicture}
\begin{axis}[width=0.94\textwidth,height=6.0cm,xlabel={Numero di alternanze $N$},ylabel={Frazione residua $L+M$},xmin=1,xmax=20,ymin=0,ymax=1,grid=major,legend columns=3,legend style={font=\footnotesize,at={(0.5,1.04)},anchor=south,draw=none},tick label style={font=\small},label style={font=\small}]
\addplot[thick,black,mark=*] coordinates {(1,0.6680496) (2,0.5248837) (3,0.4766819) (4,0.4592498) (5,0.452253) (6,0.4493265) (7,0.4481709) (8,0.4478442) (9,0.4479215) (10,0.4481971) (11,0.4485664) (12,0.4489743) (13,0.4493913) (14,0.4498015) (15,0.4501966) (16,0.4505728) (17,0.4509283) (18,0.451263) (19,0.4515773) (20,0.4518724)};
\addlegendentry{Complementari}
\addplot[thick,blue!75!black,mark=*] coordinates {(1,0.2523549) (2,0.3064317) (3,0.3328908) (4,0.3483364) (5,0.3584219) (6,0.3655152) (7,0.3707725) (8,0.3748238) (9,0.3780406) (10,0.3806565) (11,0.3828253) (12,0.3846526) (13,0.3862129) (14,0.3875609) (15,0.3887371) (16,0.3897723) (17,0.3906905) (18,0.3915104) (19,0.392247) (20,0.3929124)};
\addlegendentry{Preparazione poi estrazione}
\addplot[thick,red!75!black,mark=*] coordinates {(1,0.8798123) (2,0.8798123) (3,0.8798123) (4,0.8798123) (5,0.8798123) (6,0.8798123) (7,0.8798123) (8,0.8798123) (9,0.8798123) (10,0.8798123) (11,0.8798123) (12,0.8798123) (13,0.8798123) (14,0.8798123) (15,0.8798123) (16,0.8798123) (17,0.8798123) (18,0.8798123) (19,0.8798123) (20,0.8798123)};
\addlegendentry{Identici}
\end{axis}
\end{tikzpicture}
\caption{Modello a stati conformazionali. Ogni curva mantiene invariati il tempo totale nel liquido D e quello nel liquido T. Alternare pi\`u spesso pu\`o aiutare o peggiorare: dipende dai tassi, qui illustrativi. D e T sono stati di trattamento ipotetici, non misure di DMF e THF.}
\label{fig:cycles}
\end{figure}

\section{Il solvente deve anche portare via il materiale}
\label{sec:transport}

\subsection{Bilancio fra superficie, strato vicino e bagno}
Sia $\Gamma$ la massa accessibile per area, $I$ una frazione
schermata, $c$ la concentrazione nel bagno, $c_s$ quella prossima
alla superficie. Con $A$ area, $V$ volume, $Q$ portata di ricambio,
poniamo
\begin{align}
 J&=k_d\Gamma-k_a c_s=k_m(c_s-c),\label{eq:flux}\\
 J&=\lambda\Gamma-bc,\qquad
 \lambda=\frac{k_d}{1+k_a/k_m},\quad
 b=\frac{k_a}{1+k_a/k_m},\label{eq:effective}\\
 \dot I&=-k_iI,\quad
 \dot\Gamma=k_iI-J,\quad
 V\dot c=AJ-Qc.\label{eq:balance}
\end{align}
Qui $k_d,k_i$ hanno unit\`a s$^{-1}$; $k_a,k_m$ m/s;
$\Gamma,I$ kg/m$^2$; $c$ kg/m$^3$. Il fluido in entrata \`e
assunto privo del polimero considerato. Si verifica
\begin{equation}
 \frac{d}{dt}\{A(\Gamma+I)+Vc\}=-Qc.
\end{equation}
Il modello assume regime diluito, risposta lineare e parametri
costanti in ciascuna fase. Non rappresenta in dettaglio la forma
delle particelle n\'e una rete reticolata.

Con $h=V/A$, $\tau=\lambda t$, $x=\Gamma/\Gamma_0$,
$y=hc/\Gamma_0$, $z=I/\Gamma_0$,
\begin{equation}\label{eq:dimensionless}
 x'=-x+\beta y+\varepsilon z,\quad
 y'=x-(\beta+r)y,\quad z'=-\varepsilon z,
\end{equation}
dove $\beta=b/(h\lambda)$, $r=Q/(V\lambda)$ e
$\varepsilon=k_i/\lambda$. $\beta$ misura la tendenza a
ripartirsi tra superficie e bagno: \textbf{non \`e la saturazione
del solvente rispetto alla solubilit\`a del PMMA massivo}.

\subsection{Un limite che impedisce di promettere troppo al ricambio}
Da $c,I\ge0$ segue
\begin{equation}\label{eq:bound}
 \Gamma(t)\ge\Gamma_0e^{-\lambda t},\qquad \lambda\le k_d.
\end{equation}
Anche un bagno perfettamente pulito non supera il distacco
intrinseco nel modello. Un trasporto migliore aiuta quando limita
il riadsorbimento; pu\`o essere quasi inutile se domina una
barriera di distacco o una frazione inaccessibile.

Per un bagno senza ricambio, inizialmente pulito e senza $I$,
\begin{equation}\label{eq:static}
 x(\tau)=\frac{\beta+e^{-(1+\beta)\tau}}{1+\beta}.
\end{equation}
Compare un plateau pur essendo tutto reversibile: osservare un
plateau non prova una rete insolubile. La
Figura~\ref{fig:transport} mostra come cambia l'andamento con $r$.

\begin{figure}[htbp]
\centering
\begin{tikzpicture}
\begin{axis}[width=0.94\textwidth,height=6.0cm,xlabel={Tempo adimensionale $\tau=\lambda t$},ylabel={Frazione ancora aderente $x$},xmin=0,xmax=8,ymin=0,ymax=1,grid=major,legend style={font=\footnotesize,at={(0.98,0.98)},anchor=north east},tick label style={font=\small},label style={font=\small}]
\addplot[thick,black] coordinates {(0,1) (0.1,0.9093654) (0.2,0.83516) (0.3,0.7744058) (0.4,0.7246645) (0.5,0.6839397) (0.6,0.6505971) (0.7,0.6232985) (0.8,0.6009483) (0.9,0.5826494) (1,0.5676676) (1.1,0.5554016) (1.2,0.545359) (1.3,0.5371368) (1.4,0.530405) (1.5,0.5248935) (1.6,0.5203811) (1.7,0.5166866) (1.8,0.5136619) (1.9,0.5111854) (2,0.5091578) (2.1,0.5074978) (2.2,0.5061387) (2.3,0.5050259) (2.4,0.5041149) (2.5,0.503369) (2.6,0.5027583) (2.7,0.5022583) (2.8,0.5018489) (2.9,0.5015138) (3,0.5012394) (3.1,0.5010147) (3.2,0.5008308) (3.3,0.5006802) (3.4,0.5005569) (3.5,0.5004559) (3.6,0.5003733) (3.7,0.5003056) (3.8,0.5002502) (3.9,0.5002049) (4,0.5001677) (4.1,0.5001373) (4.2,0.5001124) (4.3,0.5000921) (4.4,0.5000754) (4.5,0.5000617) (4.6,0.5000505) (4.7,0.5000414) (4.8,0.5000339) (4.9,0.5000277) (5,0.5000227) (5.1,0.5000186) (5.2,0.5000152) (5.3,0.5000125) (5.4,0.5000102) (5.5,0.5000084) (5.6,0.5000068) (5.7,0.5000056) (5.8,0.5000046) (5.9,0.5000038) (6,0.5000031) (6.1,0.5000025) (6.2,0.5000021) (6.3,0.5000017) (6.4,0.5000014) (6.5,0.5000011) (6.6,0.5000009) (6.7,0.5000008) (6.8,0.5000006) (6.9,0.5000005) (7,0.5000004) (7.1,0.5000003) (7.2,0.5000003) (7.3,0.5000002) (7.4,0.5000002) (7.5,0.5000002) (7.6,0.5000001) (7.7,0.5000001) (7.8,0.5000001) (7.9,0.5000001) (8,0.5000001)};
\addlegendentry{$r=0$}
\addplot[thick,blue!75!black] coordinates {(0,1) (0.1,0.9093353) (0.2,0.834943) (0.3,0.773743) (0.4,0.7232401) (0.5,0.6814123) (0.6,0.6466215) (0.7,0.6175407) (0.8,0.5930942) (0.9,0.5724109) (1,0.5547847) (1.1,0.5396435) (1.2,0.5265238) (1.3,0.5150498) (1.4,0.5049173) (1.5,0.4958794) (1.6,0.4877362) (1.7,0.4803257) (1.8,0.4735164) (1.9,0.4672019) (2,0.4612957) (2.1,0.4557279) (2.2,0.4504416) (2.3,0.4453907) (2.4,0.4405378) (2.5,0.4358525) (2.6,0.4313101) (2.7,0.4268906) (2.8,0.4225776) (2.9,0.4183577) (3,0.4142203) (3.1,0.4101564) (3.2,0.4061589) (3.3,0.4022217) (3.4,0.3983401) (3.5,0.39451) (3.6,0.3907281) (3.7,0.3869917) (3.8,0.3832986) (3.9,0.3796468) (4,0.3760346) (4.1,0.3724608) (4.2,0.3689243) (4.3,0.3654239) (4.4,0.3619588) (4.5,0.3585283) (4.6,0.3551317) (4.7,0.3517685) (4.8,0.348438) (4.9,0.3451397) (5,0.3418733) (5.1,0.3386383) (5.2,0.3354343) (5.3,0.3322609) (5.4,0.3291178) (5.5,0.3260046) (5.6,0.3229211) (5.7,0.3198668) (5.8,0.3168416) (5.9,0.3138451) (6,0.3108769) (6.1,0.3079369) (6.2,0.3050248) (6.3,0.3021402) (6.4,0.299283) (6.5,0.2964528) (6.6,0.2936493) (6.7,0.2908724) (6.8,0.2881218) (6.9,0.2853972) (7,0.2826984) (7.1,0.2800251) (7.2,0.2773771) (7.3,0.2747541) (7.4,0.272156) (7.5,0.2695824) (7.6,0.2670331) (7.7,0.264508) (7.8,0.2620067) (7.9,0.2595291) (8,0.2570749)};
\addlegendentry{$r=0.2$}
\addplot[thick,red!75!black] coordinates {(0,1) (0.1,0.909078) (0.2,0.8331708) (0.3,0.7685775) (0.4,0.7126311) (0.5,0.6634017) (0.6,0.6194848) (0.7,0.5798518) (0.8,0.5437427) (0.9,0.5105898) (1,0.4799642) (1.1,0.4515367) (1.2,0.4250503) (1.3,0.4003011) (1.4,0.3771236) (1.5,0.3553811) (1.6,0.3349583) (1.7,0.3157562) (1.8,0.2976884) (1.9,0.2806782) (2,0.2646569) (2.1,0.2495622) (2.2,0.2353369) (2.3,0.2219286) (2.4,0.2092886) (2.5,0.1973715) (2.6,0.1861352) (2.7,0.1755401) (2.8,0.1655492) (2.9,0.1561277) (3,0.147243) (3.1,0.1388643) (3.2,0.1309626) (3.3,0.1235107) (3.4,0.116483) (3.5,0.1098553) (3.6,0.1036048) (3.7,0.09770993) (3.8,0.09215052) (3.9,0.08690745) (4,0.08196271) (4.1,0.07729933) (4.2,0.07290128) (4.3,0.06875348) (4.4,0.06484167) (4.5,0.06115243) (4.6,0.0576731) (4.7,0.05439173) (4.8,0.05129706) (4.9,0.04837846) (5,0.04562592) (5.1,0.04302999) (5.2,0.04058176) (5.3,0.03827282) (5.4,0.03609525) (5.5,0.03404158) (5.6,0.03210475) (5.7,0.03027812) (5.8,0.02855542) (5.9,0.02693073) (6,0.02539848) (6.1,0.02395341) (6.2,0.02259056) (6.3,0.02130525) (6.4,0.02009307) (6.5,0.01894985) (6.6,0.01787168) (6.7,0.01685486) (6.8,0.01589589) (6.9,0.01499147) (7,0.01413852) (7.1,0.0133341) (7.2,0.01257544) (7.3,0.01185995) (7.4,0.01118517) (7.5,0.01054878) (7.6,0.009948593) (7.7,0.009382559) (7.8,0.008848729) (7.9,0.008345273) (8,0.007870461)};
\addlegendentry{$r=2$}
\addplot[thick,green!45!black] coordinates {(0,1) (0.1,0.9074075) (0.2,0.8249234) (0.3,0.7501247) (0.4,0.6821311) (0.5,0.6203034) (0.6,0.5640801) (0.7,0.5129527) (0.8,0.4664595) (0.9,0.4241804) (1,0.3857334) (1.1,0.3507711) (1.2,0.3189778) (1.3,0.2900662) (1.4,0.2637751) (1.5,0.2398669) (1.6,0.2181258) (1.7,0.1983552) (1.8,0.1803766) (1.9,0.1640276) (2,0.1491604) (2.1,0.1356407) (2.2,0.1233465) (2.3,0.1121665) (2.4,0.1019999) (2.5,0.09275484) (2.6,0.08434769) (2.7,0.07670255) (2.8,0.06975036) (2.9,0.0634283) (3,0.05767926) (3.1,0.05245131) (3.2,0.0476972) (3.3,0.04337401) (3.4,0.03944266) (3.5,0.03586764) (3.6,0.03261665) (3.7,0.02966033) (3.8,0.02697197) (3.9,0.02452727) (4,0.02230416) (4.1,0.02028255) (4.2,0.01844417) (4.3,0.01677242) (4.4,0.0152522) (4.5,0.01386976) (4.6,0.01261263) (4.7,0.01146944) (4.8,0.01042987) (4.9,0.009484523) (5,0.008624861) (5.1,0.007843118) (5.2,0.007132231) (5.3,0.006485777) (5.4,0.005897916) (5.5,0.005363339) (5.6,0.004877215) (5.7,0.004435152) (5.8,0.004033157) (5.9,0.003667598) (6,0.003335173) (6.1,0.003032878) (6.2,0.002757983) (6.3,0.002508004) (6.4,0.002280683) (6.5,0.002073965) (6.6,0.001885984) (6.7,0.001715042) (6.8,0.001559593) (6.9,0.001418234) (7,0.001289688) (7.1,0.001172793) (7.2,0.001066493) (7.3,0.0009698274) (7.4,0.0008819238) (7.5,0.0008019877) (7.6,0.0007292968) (7.7,0.0006631945) (7.8,0.0006030837) (7.9,0.0005484211) (8,0.0004987131)};
\addlegendentry{$r=20$}
\addplot[black,dashed,domain=0:8,samples=80]{exp(-x)};
\addlegendentry{Estrazione ideale}
\end{axis}
\end{tikzpicture}
\caption{Modello di trasporto: $\beta=1$, con diversi ricambi $r$. Parametri scelti, non stimati sul campione. Il ricambio riduce il riadsorbimento ma non supera la velocit\`a intrinseca di distacco.}
\label{fig:transport}
\end{figure}

\subsection{Ricambi a volume e tempo totali invariati}
Nel caso senza frazione schermata, $I_0=0$, se si dividono
$V_{\rm tot}$ e $t_{\rm tot}$ in $N$ bagni freschi
uguali, mantenendo $\lambda$ invariato per ipotesi e ponendo
$B=Ak_a/(V_{\rm tot}k_d)$, il residuo \`e
\begin{equation}\label{eq:resource}
 x_N=\left[\frac{NB+e^{-(1+NB)\tau_{\rm tot}/N}}{1+NB}\right]^N,
 \quad \tau_{\rm tot}=\lambda t_{\rm tot}.
\end{equation}
Per $B>0$,
\begin{equation}\label{eq:resource-limit}
 \lim_{N\to\infty}x_N=
 \exp\!\left[-\frac{1-e^{-B\tau_{\rm tot}}}{B}\right].
\end{equation}
Per $B=0$ il risultato \`e $e^{-\tau_{\rm tot}}$ per ogni $N$.
Ad esempio, con $B=1$ e $\tau_{\rm tot}=5$, passare da uno a
quattro bagni porta il residuo da 0,5000 a 0,4104; il limite
di infiniti ricambi \`e 0,3704. Sono valori del modello.
La formula non giustifica bagni arbitrariamente piccoli in
laboratorio: cambiare geometria pu\`o cambiare $k_m$ e $\lambda$.

\subsection{Un parametro fisico verificato, con un uso delimitato}
Arai e collaboratori \cite{Arai1996} misurano
$D=2{,}93\times10^{-11}\,\mathrm{m^2/s}$ per PMMA atattico
di $M_w=9{,}52\times10^5$ in acetone a 25\,\degC.
\`E la diffusione di catene gi\`a sciolte. Per distanze
ipotetiche $\delta$ di 10, 100 e 1000\,\um,
$\delta^2/D$ vale circa 3,4 secondi, 5,7 minuti e 9,5 ore.
Questi ordini di grandezza mostrano il peso della geometria.
Non determinano la diffusione in DMF/THF, il tempo di ingresso
nel residuo, $k_d$, o uno spessore $\delta$ dai soli 500 rpm.

\section{Una spiegazione dell'asciugatura che si pu\`o smentire}

Se il liquido trattenuto sul campione all'uscita ha concentrazione
massica $c_f$ e volume $V_f$, la massa che pu\`o depositare \`e
limitata da
\begin{equation}\label{eq:mass-bound}
 M_{\rm dep}\le c_fV_f,\qquad
 \Gamma_{\rm dep}\le c_fh_f,\quad h_f=V_f/A.
\end{equation}
Per uno spessore equivalente e densit\`a $\rho$,
\begin{equation}\label{eq:units}
 t_{\rm eq}[\mathrm{nm}]\le
 \frac{0{,}001\;c_f[\mathrm{mg/L}]\;h_f[\mu\mathrm m]}
 {\rho[\mathrm{g/cm^3}]}.
\end{equation}
L'esempio $c_f=1$ mg/L, $h_f=10$\,\um\ e densit\`a
assunta 1,18 g/cm$^3$ d\`a al massimo 0,0085 nm medi.
Per spiegare 0,5 nm medi, con quello stesso film, servirebbero
almeno 59 mg/L. Non sono misure di questo campione.

Per usare il limite occorre delimitare la concentrazione
\emph{locale} del liquido trascinato: la media di un bagno grande
pu\`o sottostimarla. Il confronto va fatto su massa e area
complessive; l'evaporazione pu\`o concentrare il materiale in
isole alte. Il limite esclude soltanto il deposito proveniente
dall'ultimo film liquido, non il riadsorbimento avvenuto durante
tutto il bagno.

Questa misura decide se valga la pena ottimizzare l'asciugatura.
Se perfino il massimo materiale trascinabile \`e insufficiente,
``\`e tutto un deposito finale'' non spiega il residuo.
Se \`e sufficiente, il meccanismo rimane possibile ma non provato.

\section{Il programma sperimentale proposto}
\label{sec:experiment}

\subsection{Prima prova: separare i due passaggi IPA}
Si usano campioni gemelli dello stesso trasferimento e la stessa
temperatura del procedimento di Armando, da registrare. Ogni
numero della Tabella~\ref{tab:factorial} \`e un bagno distinto;
i bagni sostitutivi sono freschi. Volume, geometria, agitazione,
intervallo di trasferimento e getto di azoto devono essere
uguali. Non si lasciano asciugare i campioni fra due bagni.
R riproduce la sequenza riferita con freschezza dei bagni
standardizzata; non sappiamo se questo dettaglio coincida con
le prove gi\`a eseguite da Armando.

\begin{table}[htbp]
\caption{Quattro percorsi da 76 minuti. ``5+3'' significa due
bagni freschi distinti. Nessun risultato \`e ancora acquisito.}
\label{tab:factorial}\centering\small
\begin{tabular}{@{}lcccc@{}}\toprule
Gruppo & Primo & Intermedio & Secondo & Finale\\\midrule
R, riferimento &DMF 30&IPA 5+3&THF 30&IPA 5+3\\
I, cambia intermedio &DMF 30&DMF 5+3&THF 30&IPA 5+3\\
F, cambia finale &DMF 30&IPA 5+3&THF 30&THF 5+3\\
IF, cambia entrambi &DMF 30&DMF 5+3&THF 30&THF 5+3\\\bottomrule
\end{tabular}
\end{table}

Le varianti I e IF costituiscono il primo tentativo concreto di
evitare un eventuale riblocco prima dell'estrazione in THF.
Il confronto I contro R e IF contro F misura la sostituzione
intermedia nei due contesti. F contro R e IF contro I studia
il passaggio finale, che cambia anche evaporazione e forze
capillari. Nessuna di queste differenze va chiamata automaticamente
``precipitazione''.

Si raccoglie separatamente il liquido dopo DMF, dopo ciascun IPA
e dopo THF. Su aliquote dell'eluato reale, fuori dal campione,
si verifica l'eventuale comparsa di particelle durante l'aggiunta
di IPA, usando bianchi e PMMA/PPC dello stesso lotto. Si cerca
anche se il materiale torna solubile reintroducendo il liquido
di partenza. L'esito pu\`o motivare un percorso di scambio
graduale, ma non fissa un rapporto ottimale senza misure.

Un esito negativo della prova di torbidit\`a non esclude
aggregati molto piccoli o nucleazione selettiva sulla superficie.
Un bianco di solo \SiO\ non riproduce tutte le propriet\`a
del grafene: serve anche, se disponibile, un testimone grafenico
pulito sottoposto al medesimo eluato.

\subsection{Seconda prova: sfruttare un'eventuale mobilizzazione reversibile}
Se l'IPA intermedio \`e utile, oppure se si osserva riorganizzazione
seguita da esportazione nel THF, si prova la
Tabella~\ref{tab:cycles}. L'ipotesi \`e estrarre presto la
frazione mobile prima che ritorni aderente.

\begin{table}[htbp]
\caption{Confronto fra blocchi e alternanze, entrambi da 76 minuti:
30 minuti DMF, 16 IPA, 30 THF. Stesso numero di bagni freschi,
stesso liquido finale. Il riferimento R va misurato separatamente.}
\label{tab:cycles}\centering\small
\begin{tabular}{@{}p{0.14\textwidth}p{0.77\textwidth}@{}}\toprule
Percorso & Sequenza\\\midrule
Blocchi & DMF 30; quattro bagni IPA da 4 minuti;
quattro bagni THF da 7,5 minuti; azoto.\\[4pt]
Alternato & DMF 30; quattro ripetizioni di IPA 4 minuti
seguito da THF 7,5 minuti; azoto.\\\bottomrule
\end{tabular}
\end{table}

Non \`e un invito a moltiplicare indiscriminatamente i bagni.
Le due varianti hanno nove bagni ciascuna, incluso il DMF
iniziale, e terminano entrambe nel THF. Il confronto tiene
fissi il tempo complessivo in ogni liquido e il numero di
manipolazioni; l'ordine resta la variabile. Per ciascuna coppia
si usano volumi, recipienti e modalit\`a di ricambio identici.

Il materiale esportato nei singoli bagni \`e il primo esito
da confrontare, insieme alla superficie conservata. Se
l'alternanza migliora solo l'immagine finale ma non le altre
misure, non \`e ancora l'amplificazione del meccanismo proposto.
La ricerca di una durata ottimale dei cicli ha senso soltanto
dopo un segnale replicato a tempi totali uguali.

\subsection{Una previsione geometrica aggiuntiva}
Se domina una rimozione dal bordo di isole isolate e di altezza
costante, con velocit\`a $v$, allora
\begin{equation}\label{eq:edge}
 r(t)=\max(r_0-vt,0),\qquad
 \frac{M(t)}{M_0}=\left[\max\left(1-\frac{vt}{r_0},0\right)\right]^2.
\end{equation}
Il tempo di sparizione cresce come $r_0$. Una rimozione
volumetrica di primo ordine prevede invece lo stesso tempo
di riduzione frazionale per tutte le dimensioni. Si seguono
quindi anche raggio, altezza e volume delle stesse isole.
Una dipendenza dal raggio non dimostra da sola il meccanismo
di bordo: diffusione, eterogeneit\`a e coalescenza vanno distinti.

\section{Le misure che decidono la riuscita}

\subsection{Identit\`a: un controllo breve prima di inseguire altri liquidi}
Si confrontano residuo persistente, eluati e riferimenti PMMA,
PPC e loro sovrapposizione, trattati con la stessa storia.
Se il riferimento libero si dissolve ma il residuo sul grafene
resta, diventa plausibile una barriera interfaciale. Se persiste
anche il riferimento privo di grafene, occorre verificare
alterazione, reticolazione o materiale diverso. I 90\,\degC\
per due minuti, da soli, non provano carbonizzazione.

Un singolo picco carbonilico XPS non separa PMMA e PPC.
Si usano firme complete e, se la quantit\`a lo consente,
spettroscopia infrarossa locale o frammenti ToF-SIMS confrontati
con riferimenti. Le analisi distruttive richiedono gemelli
dedicati. I materiali ausiliari, come PDMS o adesivi, entrano
nei controlli soltanto se sono effettivamente usati. Il segnale
globale di silicio su \SiO\ non identifica un residuo siliconico.

Per studiare un'eventuale modifica chimica si analizzano anche
riferimenti esposti ai bagni. La perdita della firma dell'estere
pu\`o indicare trasformazione senza asportazione. La cromatografia
dimensionale del materiale estratto, se misurabile, riguarda la
frazione solubile e non esclude una rete insolubile rimasta sul
campione. Nessun singolo strumento chiude da solo il bilancio.

\subsection{Morfologia, integrit\`a e limite di rivelabilit\`a}
Si misurano copertura, volume apparente per area e firma chimica,
non soltanto rugosit\`a. Una pellicola uniforme pu\`o essere
liscia e ancora presente. Le stesse aree vengono registrate
prima e dopo, insieme ad aree casuali e campioni non
prescansionati. La punta AFM non deve diventare l'agente di
pulizia; in liquido si confrontano campioni nello stesso mezzo.

La rimozione va calcolata sull'intersezione delle regioni dove
il grafene \`e presente prima e dopo. Le regioni asportate
sono danno, non superficie pulita. Mappa Raman, fori,
delaminazioni e substrati testimone controllano l'integrit\`a.
Mobilit\`a, resistenza di contatto e drogaggio restano esiti
separati dalla quantit\`a di polimero.

Come primo obiettivo progettuale si pu\`o fissare una riduzione
del 90\% di copertura e volume residuo, con diminuzione
chimica coerente e almeno 99\% dell'area di grafene conservata.
Non sono risultati ottenuti n\'e standard universali. Per la
scommessa occorre concordare prima se basta tale miglioramento
oppure se si richiede residuo non rilevato a un limite dichiarato.
Quest'ultimo richiede calibrazione della sensibilit\`a e non
equivale ad assenza assoluta di ogni molecola.

Una prima esplorazione pu\`o usare un trasferimento suddiviso
fra condizioni; la conferma deve coinvolgere trasferimenti
indipendenti, inizialmente almeno tre e poi un numero adeguato
alla variabilit\`a osservata. I pixel e le immagini dello
stesso campione non sono repliche indipendenti. Si riportano
tutte le condizioni, incluse quelle che non migliorano.

\section{Decisione scientifica e prossimo risultato da cercare}

La nuova proposta consiste nel \emph{modificare l'ordine e la
tempistica di una sequenza gi\`a parzialmente attiva}, controllando
quando il residuo diventa estraibile e quando torna aderente.
Il risultato teorico \`e che questo effetto \`e possibile, ha
condizioni di segno e pu\`o essere distinto da una semplice
somma di tempi in solventi. Le simulazioni escludono una
promessa generale di efficacia dei cicli.

Se I migliora rispetto a R senza danno, il primo progresso
operativo \`e una sequenza DMF--THF senza IPA intermedio.
Se l'IPA \`e utile e il confronto alternato supera quello
a blocchi, il progresso \`e un'estrazione a cicli da ottimizzare
con un bilancio di massa. Se la massa trascinabile non pu\`o
spiegare il residuo, l'asciugatura perde priorit\`a. Se il
residuo ha un'identit\`a diversa, il bersaglio chimico va
corretto prima di scegliere altri reagenti.

Nessuna di queste conclusioni \`e stata ancora osservata sul
campione di Armando. Il contributo qui ottenuto \`e una teoria
con previsioni controllabili, calcoli riproducibili e due
interventi nuovi rispetto ai tentativi riferiti. Non viene
rivendicata priorit\`a scientifica della strategia o risoluzione
sperimentale del problema.

\appendix
\section{Verifica dei calcoli e riproduzione}

Il programma \texttt{model\_cleaning.py} richiede Python e NumPy.
Il comando
\begin{center}
\texttt{python model\_cleaning.py --out-dir risultati}
\end{center}
genera \texttt{risultati.json} e le coordinate dei grafici
in \texttt{figure-incorporate.tex}. I grafici di questo
documento contengono direttamente tali coordinate e non
dipendono da immagini esterne.

Sono stati controllati 540 casi del modello di trasporto
(6 valori di $\beta$, 6 di $r$, 3 di $\varepsilon$ e 5 tempi).
L'errore massimo di conservazione della matrice di transizione
\`e $7{,}8\times10^{-14}$; lo scarto dalla soluzione analitica
del bagno senza ricambio \`e $2{,}0\times10^{-14}$.
Il minimo elemento delle matrici \`e zero. Sono inoltre
verificati il limite di estrazione ideale, la catena
irreversibile $L\to M\to R$, l'invarianza quando i due
generatori coincidono e il limite dei ricambi a risorse fisse.
Un'integrazione indipendente a passi piccoli conferma
l'evoluzione anche con popolazione inizialmente schermata,
con scarto inferiore a $10^{-12}$ nei casi confrontati.

Questi controlli attestano la soluzione delle equazioni
implementate. Non attestano che le equazioni descrivano
correttamente il residuo reale. Le costanti $a,b,e$, $k_d$,
$k_a$, $k_m$ e la frazione inaccessibile restano da identificare
con esperimenti. Non \`e stata eseguita una simulazione
atomistica del sistema PMMA950K/PPC/grafene/\SiO\ nei
solventi effettivi.

\section{Derivazioni brevi dei limiti usati}

Per la formula~\eqref{eq:sign}, durante P la massa mobile
varia di
\[
 M(t)-M_0=\frac{aL_0-bM_0}{a+b}
                 (1-e^{-(a+b)t}).
\]
Durante G viene esportata la frazione
$k(1-e^{-(k+c)u})/(k+c)$ della massa mobile iniziale.
Il prodotto d\`a la differenza fra i due ordini.

Per~\eqref{eq:resource}, in ogni bagno la quantit\`a
$\beta$ diventa $NB$ e il tempo $\tau_{\rm tot}/N$.
Si applica~\eqref{eq:static} e si moltiplicano i fattori
perch\'e ogni nuovo bagno riparte pulito. Il logaritmo del
fattore \`e asintotico a
$-(1-e^{-B\tau_{\rm tot}})/(NB)$, da cui
\eqref{eq:resource-limit}. Lasciare un bagno carico non
equivale a questo azzeramento della concentrazione.

Per~\eqref{eq:counterexample}, in ciascuna fase D la massa
$L$ si moltiplica per $e^{-w/N}$; nella fase T accade
lo stesso a $M$. Sommando le quote mobilizzate nei vari
stadi si ottiene la formula. La funzione
$N(1-e^{-w/N})$ cresce strettamente: ponendo $x=w/N>0$,
la sua derivata \`e $1-(1+x)e^{-x}>0$.
Quindi in questo caso la rimozione peggiora con pi\`u cicli.


\begin{thebibliography}{99}
\bibitem{Allresist}
Allresist GmbH, \emph{AR-P 630--670: PMMA-Resists}, scheda tecnica,
serie AR-P 672, prodotto 672.045.
\href{https://www.allresist.de/wp-content/uploads/2020/03/AR-P630-670_Deutsch_Allresist_Produktinformation.pdf}{Scheda del produttore}.

\bibitem{Tyagi2022}
A. Tyagi, V. Mi\v{s}eikis, L. Martini et al.,
\emph{Ultra-clean high-mobility graphene on technologically relevant substrates},
Nanoscale \textbf{14} (2022), 2167--2176.
\href{https://doi.org/10.1039/D1NR05904A}{doi:10.1039/D1NR05904A}.
\href{https://www.rsc.org/suppdata/d1/nr/d1nr05904a/d1nr05904a1.pdf}{Supplemento consultato}.

\bibitem{Xiao2019}
Z. Xiao, Q. Wan e C. Durkan,
\emph{Cleaning Transferred Graphene for Optimization of Device Performance},
Advanced Materials Interfaces \textbf{6} (2019), 1801794.
\href{https://doi.org/10.1002/admi.201801794}{doi:10.1002/admi.201801794}.
\href{https://www.repository.cam.ac.uk/handle/1810/291256}{Manoscritto degli autori consultato}.

\bibitem{Manjkow1987}
J. Manjkow, J. S. Papanu, D. S. Soong, D. W. Hess e A. T. Bell,
\emph{An in situ study of dissolution and swelling behavior of
poly-(methyl methacrylate) thin films in solvent/nonsolvent binary mixtures},
Journal of Applied Physics \textbf{62} (1987), 682--688.
\href{https://doi.org/10.1063/1.339742}{doi:10.1063/1.339742}.
Consultato il riassunto primario, non il testo integrale.

\bibitem{Rooks2002}
M. J. Rooks et al., \emph{Low stress development of poly(methylmethacrylate)
for high aspect ratio structures}, Journal of Vacuum Science \& Technology B
\textbf{20} (2002), 2937--2941.
\href{https://doi.org/10.1116/1.1524971}{doi:10.1116/1.1524971}.
\href{https://nano.yale.edu/sites/default/files/files/pmma_develop_ipa_water.pdf}{Testo consultato}.

\bibitem{Huston2020}
K. J. Huston, C. E. Rice e R. G. Larson,
\emph{Forward Flux Sampling of Polymer Desorption Paths from a Solid Surface
into Dilute Solution}, Polymers \textbf{12} (2020), 2275.
\href{https://doi.org/10.3390/polym12102275}{doi:10.3390/polym12102275}.

\bibitem{Ayodele2022}
O. O. Ayodele, S. Pourianejad, A. Trofe, A. Prokofjevs e T. Ignatova,
\emph{Application of Soxhlet Extractor for Ultra-clean Graphene Transfer},
ACS Omega \textbf{7} (2022), 7297--7303.
\href{https://doi.org/10.1021/acsomega.1c07113}{doi:10.1021/acsomega.1c07113}.

\bibitem{Tang2025}
Z. Tang et al., \emph{A Closed-Loop Solvent Recycling Device for Polymer
Removal in Graphene Transfer Process}, Separations \textbf{12} (2025), 295.
\href{https://doi.org/10.3390/separations12110295}{doi:10.3390/separations12110295}.

\bibitem{Duan2022}
T. Duan, H. Li, R. Papadakis e K. Leifer,
\emph{Towards ballistic transport CVD graphene by controlled removal of
polymer residues}, Nanotechnology \textbf{33} (2022), 495704.
\href{https://doi.org/10.1088/1361-6528/ac8d9b}{doi:10.1088/1361-6528/ac8d9b}.

\bibitem{Merino2024}
J. P. Merino, S. Brosel-Oliu, G. Rius et al.,
\emph{Ethanol Solvation of Polymer Residues in Graphene Solution-Gated
Field Effect Transistors}, ACS Sustainable Chemistry \& Engineering
\textbf{12} (2024), 9133--9143.
\href{https://doi.org/10.1021/acssuschemeng.4c01538}{doi:10.1021/acssuschemeng.4c01538}.

\bibitem{Suk2013}
J. W. Suk et al., \emph{Enhancement of the Electrical Properties of Graphene
Grown by Chemical Vapor Deposition via Controlling the Effects of Polymer Residue},
Nano Letters \textbf{13} (2013), 1462--1467.
\href{https://doi.org/10.1021/nl304420b}{doi:10.1021/nl304420b}.

\bibitem{Schoenemann2018}
E. Sch\"onemann, A. Laschewsky e A. Rosenhahn,
\emph{Exploring the Long-Term Hydrolytic Behavior of Zwitterionic
Polymethacrylates and Polymethacrylamides}, Polymers \textbf{10} (2018), 639.
\href{https://doi.org/10.3390/polym10060639}{doi:10.3390/polym10060639}.

\bibitem{Jung2006}
J. H. Jung, M. Ree e H. Kim,
\emph{Acid- and base-catalyzed hydrolyses of aliphatic polycarbonates and
polyesters}, Catalysis Today \textbf{115} (2006), 283--287.
\href{https://doi.org/10.1016/j.cattod.2006.02.060}{doi:10.1016/j.cattod.2006.02.060}.
Consultati riassunto e passaggio metodologico accessibile.

\bibitem{Wang2017}
X. Wang et al., \emph{Direct Observation of Poly(Methyl Methacrylate)
Removal from a Graphene Surface}, Chemistry of Materials
\textbf{29} (2017), 2033--2039.
\href{https://doi.org/10.1021/acs.chemmater.6b03875}{doi:10.1021/acs.chemmater.6b03875}.

\bibitem{Schwartz2019}
J. J. Schwartz et al., \emph{Chemical Identification of Interlayer
Contaminants within van der Waals Heterostructures},
ACS Applied Materials \& Interfaces \textbf{11} (2019), 25578--25585.
\href{https://doi.org/10.1021/acsami.9b06594}{doi:10.1021/acsami.9b06594}.

\bibitem{Arai1996}
T. Arai, N. Sawatari, T. Yoshizaki, Y. Einaga e H. Yamakawa,
\emph{Excluded-Volume Effects on the Hydrodynamic Radius of Atactic and
Isotactic Oligo- and Poly(methyl methacrylate)s in Dilute Solution},
Macromolecules \textbf{29} (1996), 2309--2314.
\href{https://doi.org/10.1021/ma951274n}{doi:10.1021/ma951274n}.
Parametro impiegato dalla Tabella 3.
\end{thebibliography}
\end{document}
